\documentclass[11pt]{article}
\usepackage[utf8]{inputenc}
\usepackage{amsmath,amssymb}
\usepackage[margin=1in]{geometry}
\usepackage{hyperref}
\emergencystretch=3em

\title{Convergence in distribution is stable under small-probability modifications}
\author{Claude, with Daniel Murfet}
\date{2026-09-07}

\begin{document}
\maketitle
\sloppy

\begin{abstract}
A generic probabilistic lemma (Astra~\#13, route R1) needed for the division step of the posterior
ratio. If for every $\eta>0$ there are random variables $Y_n^\eta\Rightarrow W^\eta$ with
$\mu\{X_n\ne Y_n^\eta\}\le\eta$ eventually and $\mu'\{Z\ne W^\eta\}\le\eta$, then $X_n\Rightarrow Z$
(\texttt{tendstoInDistribution\_of\_approx}). The workhorse is the modification bound
$|\int f\circ X-\int f\circ Y|\le C\eta$ for a test function with $\operatorname{dist}(f(x),f(y))\le C$
(\texttt{abs\_integral\_comp\_sub\_le\_of\_ne}); the conclusion follows from Mathlib's
bounded-Lipschitz characterisation of weak convergence. Zero \texttt{sorry}/\texttt{axiom}.
\end{abstract}

\section{The things formalised}
{\small
\begin{verbatim}
variable {E} [PseudoEMetricSpace E] [MeasurableSpace E] [OpensMeasurableSpace E]
  [MeasurableEq E]
theorem integrable_comp_of_bounded (f : E → ℝ) (hf : Measurable f) (x₀ : E) (C)
    (hC : ∀ x y, dist (f x) (f y) ≤ C) (X : Ω → E) (hX : Measurable X) :
    Integrable (fun ω => f (X ω)) μ
theorem abs_integral_comp_sub_le_of_ne (f) (hf) (x₀) (C) (hC) (X Y : Ω → E) (hX hY) (η)
    (hη : 0 ≤ η) (hS : μ {ω | X ω ≠ Y ω} ≤ ENNReal.ofReal η) :
    |(∫ ω, f (X ω) ∂μ) - ∫ ω, f (Y ω) ∂μ| ≤ C * η
theorem tendstoInDistribution_of_approx [Nonempty E] [l.IsCountablyGenerated]
    (X : ι → Ω → E) (Z : Ω' → E) (hX : ∀ n, Measurable (X n)) (hZ : Measurable Z)
    (happrox : ∀ η, 0 < η → ∃ (Y : ι → Ω → E) (W : Ω' → E),
      (∀ n, Measurable (Y n)) ∧ Measurable W ∧
      TendstoInDistribution Y l W (fun _ => μ) μ' ∧
      (∀ᶠ n in l, μ {ω | X n ω ≠ Y n ω} ≤ ENNReal.ofReal η) ∧
      μ' {ω | Z ω ≠ W ω} ≤ ENNReal.ofReal η) :
    TendstoInDistribution X l Z (fun _ => μ) μ'
\end{verbatim}
}

\section{Mathematics}
Weak convergence of the laws is equivalent to $\int f\,d\mu_n\to\int f\,d\mu$ for every bounded
Lipschitz $f$ (Mathlib \texttt{tendsto\_iff\_forall\_lipschitz\_integral\_tendsto}, which is where
the countably generated filter enters). For such $f$ with oscillation $\le C$, changing a random
variable on an event of probability $\le\eta$ changes $\int f\circ X$ by at most $C\eta$: the
integrand difference is dominated by $C$ times the indicator of the disagreement event. Given
$\varepsilon>0$ take $\eta=\varepsilon/(3(C+1))$ and the approximants for this $\eta$; the three
terms of the triangle inequality are each $<\varepsilon/3$ eventually.

\section{Paper-facing meaning}
The posterior ratio $Z_n[\varphi]/Z_n[1]$ is a quotient of two jointly convergent normalised chart
integrals whose limiting denominator is a.s.\ positive but not bounded below. Mathlib's continuous
mapping theorem needs global continuity, which division lacks; the standard fix is to divide by
$\max(v,c)$ (continuous), which agrees with division except on $\{V_n<c\}$, an event of small
probability uniformly in $n$ by portmanteau. The present lemma converts that approximation into
convergence in distribution of the ratio itself (next unit).

\section{Lean notes}
\begin{itemize}
\item The disagreement event $\{X\ne Y\}$ is measurable via \texttt{measurableSet\_eq\_fun}, which
needs the class \texttt{MeasurableEq E} (instances: second-countable $T_2$ spaces, standard Borel
spaces); we assume it on $E$.
\item \texttt{Integrable.of\_bound} with the bound $\|f(x_0)\|+C$; the indicator bound integrates via
\texttt{integral\_indicator\_const} and \texttt{ENNReal.toReal\_le\_of\_le\_ofReal}.
\item Rewriting with \texttt{tendsto\_iff\_forall\_lipschitz\_integral\_tendsto} on the
\texttt{TendstoInDistribution.tendsto} field needs
\texttt{set\_option backward.isDefEq.respectTransparency.types false} (as in Mathlib's own proofs),
because \texttt{ProbabilityMeasure} is a subtype behind a definition; then
\texttt{ProbabilityMeasure.coe\_mk} and \texttt{integral\_map} expose the integrals over $\Omega$.
\item \texttt{field\_simp} closed the scalar identity by itself (a trailing \texttt{ring} errors with
``no goals''); \texttt{add\_le\_add (abs\_add\_le \_ \_) le\_rfl} for the right-associated triangle
inequality (\texttt{add\_le\_add\_right} adds on the left).
\end{itemize}

\end{document}
